\documentclass[11pt]{article}
\usepackage[T1]{fontenc}
\usepackage{amsmath}
\usepackage{amssymb}
\usepackage{amsthm}
\usepackage[hidelinks]{hyperref}
\newtheorem{theorem}{Theorem}[section]
\newtheorem{lemma}[theorem]{Lemma}
\newtheorem{corollary}[theorem]{Corollary}
\theoremstyle{definition}
\newtheorem{definition}[theorem]{Definition}
\begin{document}
\begin{center}
{\LARGE Natural number doubling}
\end{center}
\section{Natural number doubling}
\label{ll:doubling:definitions}
\begin{definition}
\label{ll:doubling:count}
A count is defined as \(\mathbb{N}\).
\end{definition}
\begin{definition}
\label{ll:doubling:double}
For every natural number \(n\), \(\operatorname{double}(n)\) is defined as \(n + n\).
\end{definition}
\begin{definition}
\label{ll:doubling:even}
For every natural number \(n\), \(n\) is even holds exactly when there exists a natural number \(k\) such that \(n = k + k\).
\end{definition}
\section{Parity}
\label{ll:doubling:parity}
\begin{theorem}
\label{ll:doubling:double-even}
For every natural number \(n\), \(\operatorname{double}(n)\) is even.
\end{theorem}
\begin{proof}
Use \(n\) as the witness.
The goal follows by reflexivity.
\end{proof}
\begin{theorem}
\label{ll:doubling:sum-even}
For every natural number \(n\), the sum of \(n\) and \(n\) is even.
\end{theorem}
\begin{proof}
Use \(n\) as the witness.
The goal follows by reflexivity.
\end{proof}
\begin{theorem}
\label{ll:doubling:zero-even}
\(0\) is even.
\end{theorem}
\begin{proof}
Use \(0\) as the witness.
The goal follows by reflexivity.
\end{proof}
\begin{theorem}
\label{ll:doubling:even-succ-succ}
For every natural number \(n\), if \(n\) is even, then \(\operatorname{succ}(\operatorname{succ}(n))\) is even.
\end{theorem}
\begin{proof}
Assume \(h\).
Consider the cases of \(h\).
Case \(\operatorname{existsintro}\) with \(k\), \(hk\):
\begin{quote}
Use \(\operatorname{succ}(k)\) as the witness.
Rewrite the goal using \(\rightarrow hk\), \(\leftarrow \texttt{Addition::succ-add}(k, k)\).
The goal follows by reflexivity.
\end{quote}
\end{proof}
\section{Sums}
\label{ll:doubling:sums}
\begin{lemma}
\label{ll:doubling:rotate}
For every natural number \(a\) and natural number \(b\), \(a + (b + (a + b)) = a + (a + (b + b))\).
\end{lemma}
\begin{proof}
Rewrite the goal using \(\rightarrow \texttt{Addition::add-left-comm}(b, a, b)\).
\end{proof}
\begin{theorem}
\label{ll:doubling:double-add}
For every natural number \(a\) and natural number \(b\), \(\operatorname{double}(a + b) = \operatorname{double}(a) + \operatorname{double}(b)\).
\end{theorem}
\begin{proof}
\begin{align*}
\operatorname{double}(a + b) &= a + b + (a + b) && \text{by } \operatorname{rfl} \\
 &= a + (b + (a + b)) && \text{by } \operatorname{addassoc}(a, b, a + b) \\
 &= a + (a + (b + b)) && \text{by } \texttt{Doubling::rotate}(a, b) \\
 &= a + a + (b + b) && \text{by } \operatorname{eqsymm}(\operatorname{addassoc}(a, a, b + b)) \\
 &= \operatorname{double}(a) + \operatorname{double}(b) && \text{by } \operatorname{rfl}
\end{align*}
\end{proof}
\end{document}
